\subsection*{Problems}

\textbf{2.1.} Show that a set function $\mu:\mathcal{F}\to [0,\infty]$ is a measure if and only if:
\begin{enumerate}[label=\arabic*.]
    \item $\mu$ is $\sigma$-additive.
    \item There exists a set $A\in \mathcal{F}$ such that $\mu(A)$ is finite.
\end{enumerate}

\textbf{2.2.}
\begin{enumerate}[label=(\alph*)]
    \item Let $A_1,A_2,A_3,\dots$ be a sequence of measurable sets with measure 0 in a measure space. Show that the union $\cup_{i=1}^{\infty} A_i$ has measure 0.
    \item Prove that if $B_1,B_2,B_3,\dots$ is a sequence of events with probability 1 in a probability space, then the intersection $\cap_{i=1}^{\infty} B_i$ has probability 1.
\end{enumerate}

\textbf{2.3.} We define the Borel algebra $\mathcal{B}(\mathbb{R})$ as the smallest $\sigma$-algebra on $\mathbb{R}$ that contains all open intervals. Show that $\mathcal{B}(\mathbb{R})$ can be generated by any one of the following collection of sets:
\begin{itemize}
    \item $\{(-\infty,b] : b\in \mathbb{R}\}$.
    \item $\{(-\infty,b) : b\in \mathbb{R}\}$.
    \item $\{(a,\infty) : a\in \mathbb{R}\}$.
    \item $\{[a,\infty) : a\in \mathbb{R}\}$.
    \item $\{(a,b] : a,b\in \mathbb{R}\}$.
\end{itemize}

\textbf{2.4.} (Union bound) For any $n$ events $A_1,A_2,\dots,A_n$, prove that
\[
\Pr(A_1\cup A_2\cup \cdots \cup A_n)\le \sum_{i=1}^{n} \Pr(A_i).
\]

Use the property of continuity from below to prove that, for any sequence of $\mathcal{F}$-measurable sets $A_n$, for $n\ge 1$, we have
\[
\Pr\left(\bigcup_{i=1}^{\infty} A_i\right)\le \sum_{i=1}^{\infty} \Pr(A_i).
\]

\textbf{2.5.} Consider the counting measure $\mu$ on a countable sample space $\Omega$. Construct a sequence of decreasing sets $A_1\supseteq A_2\supseteq A_3\supseteq \cdots$ such that $\cap_{i=1}^{\infty} A_i=\emptyset$, but $\lim_{i\to\infty} \mu(A_i)\neq 0$.

\textbf{2.6.} For $i=1,2,3,\dots$, let $A_i$ be a subset of a set $\Omega$. Prove that
\[
\liminf_{i\to\infty} A_i \subseteq \limsup_{i\to\infty} A_i.
\]

\textbf{2.7.} Find the limits superior and inferior of the following sequences of sets $(A_i)_{i\ge 1}$:
\begin{enumerate}[label=(\alph*)]
    \item $A_i=[1/i,\, 4+(-1)^i]$.
    \item $A_i$ is the interval $[0,1]$ if $i$ is odd, and $[1,2]$ if $i$ is even.
\end{enumerate}

\textbf{2.8.} Consider a sample space $\Omega=\{a,b,c,d,e,f\}$ consisting of $6$ elements. Let $A_1$ and $A_2$ denote the subsets $\{a,c,d\}$ and $\{c,d,f\}$, respectively. List the subsets in the $\sigma$-algebra generated by $A_1$ and $A_2$.

\textbf{2.9.} Let $A_1,A_2,\dots,A_n$ be $n$ arbitrary subsets of $\Omega$. Describe explicitly the $\sigma$-field generated by $A_1,A_2,\dots,A_n$, and provide an upper bound for the number of subsets it contains.

\textbf{2.10.} Show that the Borel algebra on $\mathbb{R}$ can be generated by countably many sets in $\mathbb{R}$.

\textbf{2.11.} Verify that the Cantor set is a Borel set.

\textbf{2.12.} Suppose $\mathcal{A}$ is a collection of subsets in $\Omega$ such that (a) $\Omega\in \mathcal{A}$, (b) $\mathcal{A}$ is closed under taking set difference, i.e., $A,B\in \mathcal{A}$ implies $B\setminus A\in \mathcal{A}$. Show that $\mathcal{A}$ is a field.
